\documentclass[12pt,a4paper]{article}

% ── Paquetes ──────────────────────────────────────────────────────────────
\usepackage[utf8]{inputenc}
\usepackage[T1]{fontenc}
\usepackage[spanish]{babel}
\usepackage{lmodern}
\usepackage{microtype}
\usepackage{geometry}
\geometry{margin=2.5cm}

\usepackage{amsmath,amssymb}
\usepackage{graphicx}
\usepackage{subcaption}
\usepackage{booktabs}
\usepackage{array}
\usepackage{tabularx}
\usepackage{booktabs}
\usepackage{xcolor}
\usepackage{hyperref}
\usepackage{float}
\usepackage{enumitem}
\usepackage{multirow}
\usepackage{listings}
\usepackage{caption}

\hypersetup{
    colorlinks=true,
    linkcolor=blue!70!black,
    urlcolor=blue!70!black,
    citecolor=green!60!black
}

\definecolor{codeblue}{RGB}{0,102,204}
\definecolor{codegray}{RGB}{128,128,128}
\definecolor{codegreen}{RGB}{0,153,0}
\definecolor{codebg}{RGB}{248,248,248}

\lstset{
    backgroundcolor=\color{codebg},
    basicstyle=\ttfamily\footnotesize,
    keywordstyle=\color{codeblue}\bfseries,
    commentstyle=\color{codegray}\itshape,
    stringstyle=\color{codegreen},
    breaklines=true,
    frame=single,
    rulecolor=\color{gray!40},
    language=Python
}

% ── Metadatos ─────────────────────────────────────────────────────────────
\title{
    \vspace{-1cm}
    \Large\textbf{Diseño de Redes Neuronales Confiables}\\[0.4cm]
    \large Justicia e Incertidumbre en la Concesión de Crédito\\[0.3cm]
    \normalsize Taller B4-T1 — Home Credit Default Risk
}
\author{
    \textbf{Grupo 7} \\[0.2cm]
    \small Instituto BME /
    Daniel Gallego Sánchez , Marco Corpa Criado , Javier	Fernández Guerra
}
\date{\today}

% ══════════════════════════════════════════════════════════════════════════
\begin{document}
\maketitle
\tableofcontents
\newpage

% ══════════════════════════════════════════════════════════════════════════
\section{Introducción y Objetivo}
\label{sec:intro}

El presente trabajo aborda el diseño, entrenamiento y auditoría de modelos
de clasificación neuronal para la concesión de créditos. El dataset utilizado
es el \textbf{Home Credit Default Risk} (Kaggle, 307\,507 solicitudes),
cuyo objetivo es predecir si un solicitante tendrá dificultades de pago
(\texttt{TARGET=1}, $\approx$8\%) o no (\texttt{TARGET=0}, $\approx$92\%).

Los objetivos específicos del trabajo son:

\begin{enumerate}[leftmargin=*, label=\textbf{\arabic*.}]
    \item \textbf{Arquitectura customizada}: implementar capas Keras con
          restricciones de dominio financiero.
    \item \textbf{Aprendizaje Justo (FAIR Loss)}: diseñar una función de
          coste que penalice la dependencia estadística entre las predicciones
          y el género del solicitante (\texttt{CODE\_GENDER}).
    \item \textbf{AutoML}: explorar la progresión de arquitecturas y
          justificar cada elección.
    \item \textbf{Incertidumbre}: cuantificar la varianza de las predicciones
          y analizar su relación con la calidad de los datos de entrada.
\end{enumerate}

% ══════════════════════════════════════════════════════════════════════════
\section{Análisis Exploratorio de Datos (EDA)}
\label{sec:eda}

\subsection{Dataset y features seleccionadas}

Se seleccionaron las 9 columnas más relevantes según el enunciado:
\texttt{TARGET}, \texttt{CODE\_GENDER}, \texttt{AMT\_INCOME\_TOTAL},
\texttt{AMT\_CREDIT}, \texttt{AMT\_ANNUITY}, \texttt{DAYS\_BIRTH},
\texttt{EXT\_SOURCE\_1/2/3}.

\subsection{Desbalanceo de clases}

\begin{figure}[H]
    \centering
    \includegraphics[width=0.85\textwidth]{figs/01_target_balance.pdf}
    \caption{Distribución de la variable objetivo. El 91.9\% de los
             solicitantes pagaron a tiempo; solo el 8.1\% presentó
             dificultades de pago.}
    \label{fig:balance}
\end{figure}

El fuerte desbalanceo (razón 11.4:1) tiene consecuencias directas en el
diseño del modelo:

\begin{itemize}
    \item La métrica de evaluación es \textbf{AUC-ROC}, no \textit{accuracy}:
          un modelo que prediga siempre clase 0 obtendría 91.9\% de
          \textit{accuracy} sin ninguna utilidad.
    \item Se utilizan \textbf{class weights} durante el entrenamiento
          (clase~1 ponderada $\approx$11.4× más que la clase~0) para que
          el gradiente no ignore la clase minoritaria.
\end{itemize}

\subsection{Valores ausentes}

\begin{figure}[H]
    \centering
    \includegraphics[width=0.7\textwidth]{figs/02_missing_values.pdf}
    \caption{Porcentaje de valores ausentes por variable. \texttt{EXT\_SOURCE\_1}
             presenta el mayor hueco (56.4\%), seguido de \texttt{EXT\_SOURCE\_3}
             (19.8\%).}
    \label{fig:missing}
\end{figure}

La distribución de nulos motiva dos decisiones de diseño:
\begin{enumerate}
    \item \textbf{Imputación por la mediana} (calculada exclusivamente
          sobre el conjunto de entrenamiento para evitar fuga de datos).
    \item Creación de \textbf{features de máscara binaria}
          (\texttt{EXT\_SOURCE\_*\_MISSING}) que indican si el valor
          original estaba ausente. Estas features actúan como señal
          predictiva adicional y son la materia prima del análisis de
          incertidumbre (Sección~\ref{sec:uncertainty}).
\end{enumerate}

\subsection{Variables financieras: asimetría y transformación}

\begin{figure}[H]
    \centering
    \includegraphics[width=\textwidth]{figs/03_log1p_transform.pdf}
    \caption{Distribuciones antes (fila superior) y después (fila inferior)
             de aplicar $\log(1+x)$. La transformación elimina la cola
             asimétrica extrema (outlier de \texttt{AMT\_INCOME\_TOTAL}
             $\approx 1.17 \times 10^8$) sin pérdida de información.}
    \label{fig:log1p}
\end{figure}

La transformación $\log(1+x)$ se aplica a las tres variables monetarias
(\texttt{AMT\_INCOME\_TOTAL}, \texttt{AMT\_CREDIT}, \texttt{AMT\_ANNUITY})
antes del escalado estándar. La justificación es doble: (1)~comprime colas
largas que dominarían el gradiente durante el entrenamiento, y (2)~linealiza
relaciones log-proporcionales habituales en datos financieros.

\subsection{Poder predictivo de las variables}

\begin{figure}[H]
    \centering
    \includegraphics[width=0.75\textwidth]{figs/04_ext_source_kde.pdf}
    \caption{Distribución de \texttt{EXT\_SOURCE\_1/2/3} por clase. Las
             distribuciones están claramente separadas: los solicitantes con
             impago tienen scores externos más bajos.}
    \label{fig:ext_kde}
\end{figure}

\begin{figure}[H]
    \centering
    \includegraphics[width=\textwidth]{figs/05_deciles_default_rate.pdf}
    \caption{Tasa de impago por decil de cada variable continua. Las tres
             fuentes externas muestran una relación \textbf{monotóna
             decreciente} fuerte: el modelo necesita capturar esta señal.
             \texttt{AMT\_CREDIT} y \texttt{AMT\_ANNUITY} muestran una
             \textbf{joroba} (máximo en deciles intermedios), lo que
             implica que una sola capa oculta no puede capturar esta
             no-linealidad.}
    \label{fig:deciles}
\end{figure}

\subsection{Disparidad de género}

\begin{figure}[H]
    \centering
    \includegraphics[width=0.45\textwidth]{figs/06_gender_default_rate.pdf}
    \caption{Tasa de impago por género. Los hombres (10.1\%) presentan
             mayor tasa que las mujeres (7.0\%). Las 4 filas con
             \texttt{CODE\_GENDER=XNA} fueron eliminadas.}
    \label{fig:gender}
\end{figure}

Esta disparidad de base motiva la introducción de la FAIR Loss
(Sección~\ref{sec:fair}): si el modelo optimiza solo AUC, aprenderá y
amplificará esta diferencia, lo que constituye discriminación estadística
en el sentido regulatorio.

\subsection{Interacciones no lineales: hexbin}

\begin{figure}[H]
    \centering
    \includegraphics[width=0.6\textwidth]{figs/07_hexbin_ext23.pdf}
    \caption{Hexbin de tasa de impago en el espacio
             \texttt{EXT\_SOURCE\_2} $\times$ \texttt{EXT\_SOURCE\_3}.
             La zona de máximo riesgo (rojo oscuro) se concentra en la
             esquina inferior izquierda, donde \textbf{ambas} fuentes son
             bajas simultáneamente. Este efecto de interacción no es
             capturables con una sola capa oculta.}
    \label{fig:hexbin}
\end{figure}

% ══════════════════════════════════════════════════════════════════════════
\section{Preprocesado y Pipeline de Datos}
\label{sec:pipeline}

\subsection{Decisiones de transformación}

\begin{table}[H]
\centering
\caption{Resumen de transformaciones aplicadas y su justificación.}
\label{tab:transforms}
\small
\begin{tabular}{@{}lll@{}}
\toprule
\textbf{Variable(s)} & \textbf{Transformación} & \textbf{Justificación} \\
\midrule
\texttt{AMT\_INCOME\_TOTAL} & $\log(1+x)$ & Cola asimétrica extrema ($\max = 1.17\times10^8$) \\
\texttt{AMT\_CREDIT} & $\log(1+x)$ & Mismo patrón; joroba en deciles \\
\texttt{AMT\_ANNUITY} & $\log(1+x)$ & Alta correlación con \texttt{AMT\_CREDIT} \\
\texttt{DAYS\_BIRTH} & $|x|/365$ & Convertir días negativos a edad en años \\
\texttt{CODE\_GENDER} & \texttt{M}$\to$0, \texttt{F}$\to$1 & Codificación binaria; XNA eliminado \\
\texttt{EXT\_SOURCE\_*} & Imputación mediana (train) & 20--56\% nulos; sin fuga \\
\texttt{EXT\_SOURCE\_*\_MISSING} & Máscara binaria nueva & Señal adicional + base para incertidumbre \\
\texttt{DEBT\_RATIO} & \texttt{AMT\_ANNUITY/AMT\_INCOME} & Feature de dominio financiero (clip a 3.0) \\
Variables continuas & \texttt{StandardScaler} (train) & Gradiente estable; no aplicado a binarias \\
\bottomrule
\end{tabular}
\end{table}

\subsection{Partición de datos}

Se realizó una partición estratificada \textbf{60/20/20} (train/val/test)
manteniendo la proporción de impagos en los tres conjuntos (8.07\%):

\begin{table}[H]
\centering
\caption{Tamaños de los conjuntos tras el split estratificado.}
\label{tab:splits}
\begin{tabular}{@{}lrrr@{}}
\toprule
\textbf{Conjunto} & \textbf{Muestras} & \textbf{TARGET=1} & \textbf{Tasa impago} \\
\midrule
Train      & 184\,504 & 14\,889 & 8.07\% \\
Validación &  61\,501 &  4\,963 & 8.07\% \\
Test       &  61\,502 &  4\,963 & 8.07\% \\
\midrule
Total      & 307\,507 & 24\,815 & 8.07\% \\
\bottomrule
\end{tabular}
\end{table}

El conjunto de \textbf{test permanece intocado} hasta la evaluación
final. La validación se utiliza exclusivamente para el \textit{early
stopping} y la selección de hiperparámetros.

% ══════════════════════════════════════════════════════════════════════════
\section{Arquitectura del Modelo}
\label{sec:architecture}

\subsection{Justificación del tipo de red}

\begin{table}[H]
\centering
\caption{Comparativa de familias de redes neuronales para este problema.}
\label{tab:net_types}
\small
\begin{tabularx}{\textwidth}{@{} l X X @{}}
\toprule
\textbf{Tipo} & \textbf{Requiere} & \textbf{Decisión} \\
\midrule
RNN/LSTM & Orden temporal entre muestras o features & \textcolor{red}{Descartado}: cada fila es un solicitante independiente \\
CNN & Estructura espacial/local entre features & \textcolor{red}{Descartado}: no hay vecindad entre columnas \\
Transformer & Secuencias con atención & \textcolor{red}{Descartado}: sobreingeniería para 12 features tabulares \\
\textbf{MLP denso} & -- & \textcolor{green!60!black}{\textbf{Seleccionado}}: estándar empírico para datos tabulares \\
\bottomrule
\end{tabularx}
\end{table}

\subsection{Función de activación y función de coste}

\begin{itemize}
    \item \textbf{Capas ocultas: ReLU}. Evita el problema de
          \textit{vanishing gradient} (frente a sigmoid/tanh),
          computacionalmente eficiente y empíricamente superior en
          clasificación tabular.
    \item \textbf{Capa de salida: sigmoid}. Produce una salida en $[0,1]$
          interpretable directamente como $P(\text{impago})$.
    \item \textbf{Función de coste: Binary Cross-Entropy (BCE)} con
          \textit{class weights}. BCE es la pérdida natural para
          distribuciones de Bernoulli. Los class weights compensan el
          desbalanceo sin alterar la probabilidad de las predicciones.
\end{itemize}

\subsection{Capas custom}

\subsubsection{DebtRatioLayer}

Aplica una saturación sigmoide aprendible al ratio de endeudamiento:

\begin{equation}
    f(x) = \sigma\bigl(\alpha \cdot (x - \theta)\bigr), \quad
    \alpha > 0, \quad \theta \in \mathbb{R}
    \label{eq:debt_layer}
\end{equation}

donde $\alpha$ controla la pendiente de la saturación (inicializado en 10)
y $\theta$ es el umbral de inflexión (inicializado en 0.35, referencia
regulatoria del 35\%). Ambos son \textbf{parámetros aprendibles}.

La justificación no es predictiva sino \textbf{regulatoria}: un ratio de
0.8 y uno de 1.5 son ambos ``imposible pagar''; la sigmoide garantiza que
el modelo los trate de forma similar, incrustando esta restricción de dominio
en la arquitectura.

\subsubsection{ExtSourceLayer}

Combina las tres fuentes externas en un índice de confianza ponderado:

\begin{equation}
    g(\mathbf{s}) = \sigma(w_1 s_1 + w_2 s_2 + w_3 s_3 + b),
    \quad w_i \geq 0
    \label{eq:ext_layer}
\end{equation}

Los pesos $w_i$ están restringidos a ser no negativos (la relación con el
impago es monotóna decreciente, confirmada por los deciles del EDA).
Esto hace el modelo \textbf{auditable}: tras el entrenamiento, los pesos
aprendidos revelan la importancia relativa que el modelo asignó a cada
fuente externa.

\subsection{Progresión de modelos}

\begin{table}[H]
\centering
\caption{Comparativa de modelos entrenados (AUC en test, clase 1 ponderada 11.4×).}
\label{tab:models}
\begin{tabular}{@{}llrrr@{}}
\toprule
\textbf{Modelo} & \textbf{Arquitectura} & \textbf{AUC val} & \textbf{AUC test} & \textbf{Paráms} \\
\midrule
M0 & Reg. Logística (0 capas ocultas) & 0.7275 & 0.7335 & 13 \\
M1 & MLP (Dense 64, ReLU)             & 0.7416 & 0.7449 & 897 \\
M2 & MLP (Dense 128→64, ReLU)         & 0.7411 & 0.7447 & 9\,985 \\
M3 & M2 + Dropout(0.3, 0.2)           & 0.7424 & 0.7452 & 9\,985 \\
M4 & M3 + DebtRatioLayer              & 0.7424 & 0.7449 & 9\,988 \\
M5 & M4 + ReduceLROnPlateau           & 0.7428 & 0.7449 & 9\,988 \\
\textbf{M6} & M5 + ExtSourceLayer (Dual Custom) & \textbf{0.7421} & \textbf{0.7451} & \textbf{9\,991} \\
\bottomrule
\end{tabular}
\end{table}

\begin{figure}[H]
    \centering
    \begin{subfigure}[b]{0.48\textwidth}
        \includegraphics[width=\textwidth]{figs/curves_M0.pdf}
        \caption{M0 — Regresión Logística}
    \end{subfigure}
    \hfill
    \begin{subfigure}[b]{0.48\textwidth}
        \includegraphics[width=\textwidth]{figs/curves_M3.pdf}
        \caption{M3 — MLP + Dropout}
    \end{subfigure}
    \\[0.5cm]
    \begin{subfigure}[b]{0.48\textwidth}
        \includegraphics[width=\textwidth]{figs/curves_M6.pdf}
        \caption{M6 — Dual Custom + Dropout}
    \end{subfigure}
    \caption{Curvas de entrenamiento y validación (selección representativa).
             El modelo base (M0) converge más lento y a menor AUC. M6 muestra
             convergencia suave con \textit{early stopping} efectivo.}
    \label{fig:curves}
\end{figure}

\begin{figure}[H]
    \centering
    \includegraphics[width=0.6\textwidth]{figs/arch_M6_DualCustom.pdf}
    \caption{Arquitectura del modelo M6 (Dual Custom + Dropout). Las dos ramas
             paralelas (DebtRatioLayer y ExtSourceLayer) se concatenan con la
             salida de las capas densas antes de la neurona de clasificación.}
    \label{fig:arch_m6}
\end{figure}

\paragraph{Justificación de la profundidad.}
El análisis de deciles (Figura~\ref{fig:deciles}) reveló que
\texttt{AMT\_CREDIT} y \texttt{AMT\_ANNUITY} presentan una relación en
forma de \textbf{joroba} con la tasa de impago: los créditos medios tienen
mayor tasa que los muy bajos o los muy altos. Una sola capa ReLU
(lineal por partes con una única inflexión) no puede capturar este máximo
interior; se requieren al menos dos capas ocultas.

El hexbin (Figura~\ref{fig:hexbin}) muestra además un \textbf{efecto de
interacción} entre \texttt{EXT\_SOURCE\_2} y \texttt{EXT\_SOURCE\_3} que
tampoco es modelable linealmente.

% ══════════════════════════════════════════════════════════════════════════
\section{Aprendizaje Justo — FAIR Loss}
\label{sec:fair}

\subsection{Formulación}

La función de coste customizada combina el error de clasificación con
una penalización por \textbf{Paridad Demográfica} (DP):

\begin{equation}
    \mathcal{L}_{\text{fair}} =
    \underbrace{\text{BCE}(y, \hat{y})}_{\text{clasificación}}
    + \lambda \cdot
    \underbrace{\bigl|\mathbb{E}[\hat{y} \mid G=M] -
                       \mathbb{E}[\hat{y} \mid G=F]\bigr|}_{\text{penalización DP}}
    \label{eq:fair_loss}
\end{equation}

\subsection{Justificación de la métrica de equidad}

\begin{table}[H]
\centering
\caption{Comparativa de métricas de equidad consideradas.}
\label{tab:fairness_metrics}
\small
\begin{tabular}{@{}llll@{}}
\toprule
\textbf{Métrica} & \textbf{Diferenciable} & \textbf{Requiere $y_{\text{true}}$} & \textbf{Decisión} \\
\midrule
Paridad Demográfica (DP) & \checkmark & No & \textbf{Seleccionada} \\
Equalized Odds           & Parcial    & Sí & Descartada (complica el grafo) \\
Mutual Information       & \texttimes & No & Descartada (no diferenciable) \\
Correlación $|\text{Corr}(\hat{y},G)|$ & \checkmark & No & Descartada (solo lineal) \\
\bottomrule
\end{tabular}
\end{table}

La DP mide si la tasa de aprobación predicha es similar entre grupos
demográficos, que es el requisito regulatorio explícito en la Directiva
Europea de IA y la ECOA (Equal Credit Opportunity Act).

\subsection{Curva de Pareto}

\begin{figure}[H]
    \centering
    \includegraphics[width=\textwidth]{figs/08_pareto_fair.pdf}
    \caption{Izquierda: curva de Pareto entre AUC (precisión) y DP Gap
             (discriminación). Cada punto corresponde a un valor de $\lambda$.
             Derecha: evolución de la predicción media por género al aumentar
             $\lambda$. Las líneas convergen hacia un valor similar, indicando
             que el modelo aprende a no discriminar.}
    \label{fig:pareto}
\end{figure}

\begin{table}[H]
\centering
\caption{Resultados del experimento Pareto (evaluación en test).}
\label{tab:pareto}
\begin{tabular}{@{}rrrrrr@{}}
\toprule
$\lambda$ & AUC test & DP Gap & $\hat{p}(M)$ & $\hat{p}(F)$ & $\Delta$AUC \\
\midrule
0.0 & 0.7463 & 0.0300 & 0.0998 & 0.0697 & 0.0000 \\
0.1 & 0.7440 & 0.0079 & 0.0842 & 0.0763 & $-$0.0023 \\
\textbf{0.5} & \textbf{0.7429} & \textbf{0.0034} & \textbf{0.0800} & \textbf{0.0766} & $-$\textbf{0.0034} \\
1.0 & 0.7425 & 0.0044 & 0.0850 & 0.0805 & $-$0.0038 \\
2.0 & 0.7402 & 0.0023 & 0.0756 & 0.0732 & $-$0.0060 \\
5.0 & 0.7352 & 0.0019 & 0.0733 & 0.0714 & $-$0.0110 \\
\bottomrule
\end{tabular}
\end{table}

\paragraph{Modelo recomendado: $\lambda = 0.5$.}
De $\lambda=0$ a $\lambda=0.5$ el DP Gap se reduce un \textbf{89\%}
(de 0.0300 a 0.0034) sacrificando únicamente \textbf{0.0034 puntos de AUC}.
A partir de $\lambda=0.5$ se produce un rendimiento decreciente: la ganancia
en equidad es mínima mientras el coste en AUC se acelera.

% ══════════════════════════════════════════════════════════════════════════
\section{Incertidumbre — MC Dropout}
\label{sec:uncertainty}

\subsection{Método seleccionado}

Se utiliza \textbf{Monte Carlo Dropout} (Gal \& Ghahramani, 2016), que
mantiene activas las capas Dropout durante la inferencia y realiza $N=100$
pasadas hacia adelante por muestra. La varianza entre pasadas mide la
incertidumbre epistémica del modelo:

\begin{equation}
    \hat{\sigma}^2_i = \frac{1}{N}\sum_{t=1}^{N}(\hat{y}_i^{(t)} - \bar{y}_i)^2,
    \quad \bar{y}_i = \frac{1}{N}\sum_{t=1}^{N}\hat{y}_i^{(t)}
    \label{eq:mc_var}
\end{equation}

La elección de MC Dropout frente a alternativas (Redes Bayesianas,
Deep Ensembles) se justifica porque M6 ya incorpora capas Dropout: no
requiere rediseño arquitectónico y su fundamento teórico como
aproximación variacional bayesiana está bien establecido.

\subsection{Distribución de incertidumbre por clase}

\begin{figure}[H]
    \centering
    \includegraphics[width=\textwidth]{figs/09_uncertainty_by_class.pdf}
    \caption{Distribución de la varianza MC Dropout por clase TARGET.
             M6 base (izq.) muestra distribuciones casi idénticas entre
             clases: no reconoce que los impagos son más difíciles de
             predecir. FAIR $\lambda=0.5$ (dcha.) presenta una cola más
             larga para TARGET=1, indicando mayor incertidumbre en la
             clase de riesgo.}
    \label{fig:unc_class}
\end{figure}

\subsection{Incertidumbre vs calidad de EXT\_SOURCE}

\begin{figure}[H]
    \centering
    \includegraphics[width=\textwidth]{figs/10_uncertainty_vs_missing.pdf}
    \caption{Varianza media en función del número de fuentes externas
             ausentes (imputadas con la mediana). El modelo FAIR muestra
             un patrón \textbf{monotóno creciente} (hasta $+73\%$ con
             las tres fuentes imputadas), mientras que M6 base es
             prácticamente plano.}
    \label{fig:unc_missing}
\end{figure}

\subsection{Resultados}

\begin{table}[H]
\centering
\caption{Resumen de métricas de incertidumbre (test, $N=100$ pasadas MC Dropout).}
\label{tab:uncertainty}
\begin{tabular}{@{}lrr@{}}
\toprule
\textbf{Métrica} & \textbf{M6 base} & \textbf{FAIR $\lambda=0.5$} \\
\midrule
AUC predicción puntual              & 0.7451 & 0.7429 \\
AUC media MC Dropout                & \textbf{0.7455} & \textbf{0.7433} \\
\midrule
Varianza media — TARGET=0 (pagó)    & 0.001422 & 0.000147 \\
Varianza media — TARGET=1 (impago)  & 0.001306 & \textbf{0.000364} \\
Ratio var(T=1)/var(T=0)             & 0.92 ($-8\%$) & \textbf{2.48 ($+148\%$)} \\
\midrule
Varianza — 0 EXT\_SOURCE ausentes   & 0.001472 & 0.000163 \\
Varianza — 1 EXT\_SOURCE ausente    & 0.001354 & 0.000161 \\
Varianza — 2 EXT\_SOURCE ausentes   & 0.001489 & 0.000185 \\
Varianza — 3 EXT\_SOURCE ausentes   & 0.001414 & \textbf{0.000282} \\
\bottomrule
\end{tabular}
\end{table}

\paragraph{Hallazgos clave.}
\begin{enumerate}
    \item \textbf{M6 base no distingue incertidumbre por clase}: la varianza
          en impagos es un 8\% \emph{menor} que en buenos pagadores,
          comportamiento opuesto al deseable.
    \item \textbf{FAIR $\lambda=0.5$ muestra calibración superior}: la
          varianza en impagos es 2.5× mayor que en buenos pagadores. El
          modelo reconoce que la clase minoritaria es más difícil.
    \item \textbf{FAIR detecta degradación por datos imputados}: la varianza
          crece monotónamente con el número de fuentes externas ausentes
          (+73\% con las tres imputadas).
    \item \textbf{MC-mean mejora el AUC}: la media de 100 pasadas supera
          ligeramente a la predicción puntual en ambos modelos (+0.0004),
          confirmando el efecto ensemble del MC Dropout.
\end{enumerate}

La restricción de fairness, además de reducir la discriminación demográfica,
produce un modelo \textbf{más honesto respecto a sus limitaciones de
información}: es más incierto cuando no tiene datos de calidad.

% ══════════════════════════════════════════════════════════════════════════
\section{Conclusiones}
\label{sec:conclusions}

\begin{table}[H]
\centering
\caption{Comparativa final: modelo base vs mejor modelo FAIR (test).}
\label{tab:final}
\begin{tabular}{@{}lrr@{}}
\toprule
\textbf{Métrica} & \textbf{M6 base} & \textbf{FAIR $\lambda=0.5$} \\
\midrule
AUC-ROC (test)     & \textbf{0.7451} & 0.7429 \\
DP Gap             & 0.0300          & \textbf{0.0034} \\
$\hat{p}$(impago $|$ H) & 0.0998    & \textbf{0.0800} \\
$\hat{p}$(impago $|$ M) & 0.0697    & \textbf{0.0766} \\
Varianza ratio T1/T0 & 0.92 (peor) & \textbf{2.48 (mejor)} \\
\bottomrule
\end{tabular}
\end{table}

El sistema desarrollado alcanza un AUC de 0.745 con el subconjunto de 12
features seleccionadas. La introducción de la FAIR Loss con $\lambda=0.5$
reduce la discriminación demográfica un 89\% (de 0.030 a 0.003 de DP Gap)
con un coste de solo 0.003 puntos de AUC, un tradeoff favorable desde el
punto de vista regulatorio.

El análisis de incertidumbre mediante MC Dropout revela que el modelo FAIR
es intrínsecamente más honesto: reconoce cuándo trabaja con información de
baja calidad (fuentes externas imputadas) y cuándo enfrenta la clase de
mayor riesgo (impago), aumentando su varianza de predicción en ambos casos.

\end{document}